%% ma: L10-B13-0105
%% lop: 10
%% chuong: 5
%% bai: 13
%% dang: 10.13.13
%% loai: DS
%% muc: [NB, TH, TH, VD]
%% dap_an: "ĐSSĐ"
%% nguon: "chuyen-de-cac-so-dac-trung-cua-mau-so-lieu-khong-ghep-nhom-toan-10-chuong-trinh-moi.pdf (GV. Phan Nhật Linh) (P03), Bài 02 Dạng 2 Phần II Câu 9"
%% kien_thuc: "Lớp 10 Bài 13 (số trung bình, trung vị, tứ phân vị, mốt)"
%% trang_thai: chinh-thuc
%% ghi_chu: "Lời giải gốc ý b) tính ra 19,93 (sai, đúng là 22,175); kết luận Sai không đổi. Tự duyệt theo yêu cầu giáo viên 22/10/2026."
\begin{ex}
Sản lượng lúa (tạ) của $40$ thửa ruộng thí nghiệm có cùng diện tích được trình bày trong bảng tần số sau:
\begin{center}
\begin{tabular}{|c|c|c|c|c|c|}\hline
Sản lượng & $20$ & $21$ & $22$ & $23$ & $24$\\\hline
Tần số & $5$ & $8$ & $n$ & $m$ & $6$\\\hline
\end{tabular}
\end{center}
Xét tính đúng sai của các khẳng định sau:
\choiceTF{\True Với $m=13$, $n=8$ thì mốt của mẫu số liệu là $23$ tạ}{Với $m=13$, $n=8$ thì sản lượng trung bình của $40$ thửa ruộng, làm tròn đến hàng phần trăm, là $19{,}92$ tạ}{Với $m=13$, $n=8$ thì trung vị của mẫu số liệu là $21$}{\True Nếu sản lượng trung bình của $40$ thửa ruộng là $22{,}1$ tạ thì $n=11$}
\loigiai{a) Đúng: với $n=8$, $m=13$, giá trị $23$ có tần số lớn nhất.

b) Sai: $\overline x=\dfrac{5\cdot20+8\cdot21+8\cdot22+13\cdot23+6\cdot24}{40}=\dfrac{887}{40}\approx22{,}18$.

c) Sai: tần số tích lũy $5;13;21$ nên $x_{20}=x_{21}=22$, $M_e=22$.

d) Đúng: $n+m=21$ và $22n+23m=472$ nên $n=11$, $m=10$.}
\end{ex}
